
\section{萨德--斯梅尔定理}\label{sec:A-5}

Smale~\srcite{381} 证明了下面这个\BookRecordIndex{T0296}萨德定理的无限维版本。

\begin{theorem}[萨德--斯梅尔]\label{thm:A-5-1}
设 $X,Y$ 是\BookRecordIndex{T0213}可分巴拿赫空间，$U\subset X$ 是开集。
设 $f:U\to Y$ 是 $C^\ell$ 类弗雷德霍姆映射，其中
\[
  \ell\ge\max\{1,\operatorname{index}(f)+1\}.
\]
则 $f$ 的\BookRecordIndex{T0409}正则值集
\[
  Y_{\mathrm{reg}}(f):=
  \bigl\{y\in Y\mid x\in U,\ f(x)=y
  \Longrightarrow\operatorname{im}df(x)=Y\bigr\}
\]
在贝尔意义下是\BookRecordIndex{T0300}剩余集，即它包含可数个开稠密集的交。
\end{theorem}

\BookRecordIndex{T0213}可分性条件不可或缺。若巴拿赫空间 $X$ 有稠密序列，就称 $X$ \emph{\BookRecordIndex{T0213}可分}。
每个度量空间都是仿紧的，因此每个巴拿赫空间也仿紧。这意味着 $X$ 的每个开覆盖
都可取局部有限加细。在\BookRecordIndex{T0213}可分情形下，每个局部有限覆盖都是可数的。
\footnote{若 $\{U_\alpha\}_\alpha$ 是局部有限覆盖，$\{x_i\}_i$ 是稠密序列，
则满足 $x_i\in U_\alpha$ 的对 $(\alpha,i)$ 只有可数多个。
由于每个 $U_\alpha$ 都含某个 $x_i$，映射 $(\alpha,i)\mapsto\alpha$ 是满射。}
既然可数加细的存在意味着可数子覆盖的存在，就得到以下结论。

\begin{proposition}\label{prop:A-5-2}
设 $X$ 是\BookRecordIndex{T0213}可分巴拿赫空间。则 $X$ 的每个开覆盖都有可数子覆盖。
\end{proposition}

\begin{proof}[定理~\ref{thm:A-5-1} 的证明]
我们证明：每个 $x_0\in U$ 都有开邻域 $U_0$，使得对每个闭子集
$V_0\subset U_0$，$f$ 在 $V_0$ 上的\BookRecordIndex{T0409}正则值集均为 $Y$ 中的开稠密集。
不失一般性，在定义域和陪域中分别平移坐标，使 $x_0=0\in U$ 且
$f(0)=0$，并取 $D:=df(0)$ 的\BookRecordIndex{T0252}拟逆 $T:Y\to X$。
\sourcefootnote{原书此处只令 $x_0=0$，未说明还须将靶坐标平移
$-f(x_0)$；所引用的定理~\ref{thm:A-4-3} 以 $f(0)=0$ 为假设。
以下 $f$ 与 $y$ 均在平移后的坐标中理解。}
由定理~\ref{thm:A-4-3}，存在零点的开邻域 $W\subset X$、
映到开集 $g(W)\subset X$ 上的 $C^\ell$ 微分同胚
$g:W\to g(W)$，以及 $C^\ell$ 映射 $f_0:W\to\ker T$，使
\[
  f\circ g=f_0+D.
\]
必要时缩小 $W$，可令它有界。
\sourcefootnote{原书随后用到“$V$ 有界”，但此前只要求 $V\subset W$；
这里把局部邻域 $W$ 缩为有界集，使这一推理成立。}
我们断言 $U_0:=g(W)$ 就是所需邻域。为此，回忆存在分解
\[
  X=X_0\oplus X_1,\qquad Y=Y_0\oplus Y_1,
\]
其中
\[
  X_0:=\ker D,\qquad X_1:=\operatorname{im}T,\qquad
  Y_0:=\ker T,\qquad Y_1:=\operatorname{im}D.
\]
把 $g:W\to X$ 看作 $X$ 上的坐标卡，将方程 $f(g(x))=y$ 写成
\[
  y_0=f_0(x_0,x_1),\qquad y_1=D_1x_1,
\]
其中 $x=x_0+x_1\in W$、$x_i\in X_i$，而 $D_1:X_1\to Y_1$
表示 $D$ 在 $X_1$ 上的限制。由此可知，$y=y_0+y_1$ 是
$f|_{g(W)}$ 的\BookRecordIndex{T0409}正则值，当且仅当 $y_0$ 是有限维映射
\[
  \bigl\{x_0\in X_0\mid x_0+D_1^{-1}y_1\in W\bigr\}
  \longrightarrow Y_0:
  x_0\longmapsto f_0(x_0,D_1^{-1}y_1)
\]
的\BookRecordIndex{T0409}正则值。

由于 $\dim X_0-\dim Y_0=\operatorname{index}(f)$，且
$\ell\ge\max\{1,\operatorname{index}(f)+1\}$，
有限维向量空间之间 $C^\ell$ 映射的\BookRecordIndex{T0296}萨德定理
（光滑映射见~\srcite{289}，$C^\ell$ 映射见~\srcite{2}）
表明 $f|_{g(W)}$ 的\BookRecordIndex{T0409}正则值集
\[
  Y_{\mathrm{reg}}(f;g(W)):=
  \bigl\{y\in Y\mid x\in g(W),\ f(x)=y
  \Longrightarrow\operatorname{im}df(x)=Y\bigr\}
\]
在 $Y$ 中稠密。因此，对每个子集 $V\subset W$，
$Y_{\mathrm{reg}}(f;g(V))\supset Y_{\mathrm{reg}}(f;g(W))$ 也是稠密的。

下面证明：若 $V\subset W$ 是 $X$ 中的闭集，则
$Y_{\mathrm{reg}}(f;g(V))$ 是开集。设
\[
  y_\nu=y_{\nu,0}+y_{\nu,1}\in Y
\]
是一列趋于 $y$ 的 $f|_{g(V)}$ 奇异值。于是存在 $x_\nu\in V$，使
$f(g(x_\nu))=y_\nu$，且 $df(g(x_\nu))$ 不满射。序列
\[
  x_{\nu,1}:=TDx_\nu=Ty_\nu
\]
收敛；又因为 $V$ 有界，序列
\[
  x_{\nu,0}:=x_\nu-TDx_\nu\in\ker D
\]
有界。$\ker D$ 有限维，因此可取子序列，使 $x_{\nu,0}$ 也收敛，
从而 $x_\nu=x_{\nu,0}+x_{\nu,1}$ 收敛。
由于 $V$ 是闭集，其极限 $x:=\lim_{\nu\to\infty}x_\nu$ 仍属于 $V$，
且 $f(g(x))=y$。此外，$df(g(x))$ 是具有非平凡余核的算子列的极限，
因此也不可能满射。故 $y$ 是 $f|_{g(V)}$ 的奇异值。

这样，每个 $x\in U$ 都有一个开邻域 $U_0$，使得对其每个闭子集
$V_0\subset U_0$，集合 $Y_{\mathrm{reg}}(f;V_0)$ 在 $Y$ 中开且稠密。
由于 $X$ \BookRecordIndex{T0213}可分，由命题~\ref{prop:A-5-2}，开集 $U$ 可由可数多个
这样的开邻域 $U_i$ 覆盖。因此，$U$ 也可由可数多个闭集 $V_i$ 覆盖，
且对每个 $i$，$Y_{\mathrm{reg}}(f;V_i)$ 都在 $Y$ 中开且稠密。
于是
\[
  Y_{\mathrm{reg}}(f)=\bigcap_i Y_{\mathrm{reg}}(f;V_i)
\]
是 $Y$ 中可数个开稠密集的交。
\sourcefootnote{原书从“对每个闭子集 $V_0\subset U_0$”直接过渡到最后的闭覆盖。
前面的闭性证明实际要求坐标逆像 $g^{-1}(V_0)$ 在整个 $X$ 中闭且有界；
仅知 $V_0$ 在 $U_0$ 中相对闭并不足够。一个可行的覆盖步骤是：
在各 $W$ 中选闭球 $B\subset W$，由\BookRecordIndex{T0213}可分性选出使 $g(B)$ 覆盖 $U$ 的可数族，
再直接对这些坐标闭球应用前面的闭性论证；原书未交代这一选择。}
\end{proof}
